\begin{minipage}{0.78\linewidth}
    Un pendule simple est constitué d'une boule de masse $m=\qty{50}{\g}$ accrochée
    au bout d'un fil de longueur $\ell=\qty{30}{\cm}$, de masse négligeable. La
    boule reçoit en $A$ une impulsion qui fait remonter jusqu'en $B$, de telle
    manière que le pendule fait alors un angle $\alpha=\ang{30}$ avec la la
    verticale.
    \begin{enumerate}
        \item Calculer le travail du poids de la boule entre $A$ et
              $B$.~\textbf{(1pt)}
        \item Quel serait le travail du poids de la boule, si le pendule faisait un
              tour complet~? Justifier.~\textbf{(0,5pt)}
    \end{enumerate}
\end{minipage}
\hfill
\begin{minipage}{0.20\linewidth}
      \begin{figure}[H]
          \centering
          \begin{tikzpicture}
              \def\th{30}
              \newdimen\el \el=3cm
              \node[draw,thick,pattern=north east lines,
                    minimum width=1cm,minimum height=3mm] (S) {};
              \coordinate (O) at (S.south);
              \path[draw,fill=DarkTurquoise]
                  (O) -- ([yshift=-1cm]O) arc (270:270+\th:1cm) -- cycle;
              \node at (270+\th/2:1.5cm) {$\alpha$};
              \node[circle,fill,inner sep=1.1pt,
                    label={[anchor=north east,inner sep=1.8pt]:$O$}]
                    at (O) {};
              \draw[densely dashed]
                  (O) +(270:\el) arc (270:270+\th:\el);
              \foreach \p/\a in {A/0,B/\th}{%
                  \node[circle,ball color=black!60,inner sep=0pt,
                        minimum size=4mm,
                        label={[inner sep=1.6pt]-90+\a:$\p$}]
                        (\p) at ($(O)+(270+\a:\el)$) {};
                  \fill (\p.center) circle (1pt);
                  \draw[thick] (O) -- (\p);
              }
              \node[anchor=180+\th,inner sep=1.4pt]
                  at ($(O)!.5!(B.center)$) {$\ell$};
          \end{tikzpicture}
      \end{figure}
\end{minipage}

\begin{donnees}
    ${g=\qty{10}{\N\per\kg}}$.
\end{donnees}
